\chapter{Escrita Acadêmica}

A escrita acadêmica deve ser clara, objetiva, precisa e responsável. Seu objetivo é comunicar uma investigação, não impressionar pelo excesso de termos difíceis. Um texto acadêmico de qualidade permite que o leitor compreenda o problema, acompanhe o método, avalie os argumentos e verifique as fontes.

\section{Clareza e precisão}

Clareza não significa simplificação excessiva. Significa escolher palavras adequadas, organizar ideias em ordem lógica e evitar ambiguidades. Precisão significa usar conceitos de modo consistente, definir termos importantes e evitar generalizações sem evidência.

Evite expressões vagas, como ``muitos autores dizem'', ``a sociedade precisa'', ``desde sempre'', ``é evidente'' ou ``todos sabem''. Afirmações acadêmicas precisam de delimitação, fonte ou argumento.

\section{Coesão e progressão}

O texto deve avançar. Cada parágrafo precisa contribuir para a construção do argumento. Um parágrafo acadêmico geralmente apresenta uma ideia central, desenvolve essa ideia e se conecta ao parágrafo seguinte.

Problemas comuns de progressão incluem:

\begin{itemize}
  \item parágrafos que apenas acumulam citações;
  \item repetição da mesma ideia com palavras diferentes;
  \item mudança brusca de assunto;
  \item conceitos usados antes de serem definidos;
  \item resultados apresentados antes do método;
  \item conclusão que não responde aos objetivos.
\end{itemize}

\section{Voz autoral}

Ter voz autoral não significa escrever em tom pessoal ou opinativo. Significa conduzir o texto, selecionar fontes, estabelecer relações, interpretar dados e assumir responsabilidade pelas conclusões.

O estudante deve evitar dois extremos: desaparecer atrás de citações ou escrever sem diálogo com a literatura. Um bom TCC mostra que o autor leu, compreendeu e usou fontes para construir análise própria.

\section{Uso de primeira pessoa}

A aceitação de primeira pessoa varia por área e por orientação institucional. Em muitos trabalhos, recomenda-se linguagem impessoal ou primeira pessoa do plural de modo controlado. O mais importante é manter consistência e sobriedade.

Independentemente da escolha, o texto deve evitar informalidade, subjetivismo excessivo e frases que fragilizem a análise, como ``eu acho'', ``na minha opinião'' ou ``acredito que'', quando não houver base argumentativa.

\section{Revisão textual}

Revisar não é apenas corrigir ortografia. A revisão deve observar coerência, ordem das seções, repetição, citações, referências, tabelas, figuras, títulos, paginação e conformidade com o padrão institucional.

Uma estratégia eficiente é revisar em camadas:

\begin{enumerate}
  \item coerência entre problema, objetivos, método e conclusão;
  \item estrutura das seções e progressão argumentativa;
  \item parágrafos, frases e transições;
  \item citações, referências, notas e fontes;
  \item normalização, sumário, listas e paginação;
  \item leitura final do PDF.
\end{enumerate}

\caixaorientacao{Sempre que possível, deixe intervalo entre escrita e revisão. O estudante identifica mais problemas quando relê o texto depois de alguma distância.}
